\documentclass[letterpaper,11pt]{article}

% ---------- Pacotes ----------
\usepackage[utf8]{inputenc}
\usepackage[T1]{fontenc}
\usepackage{mathptmx}                % fonte estilo Times New Roman
\usepackage[top=1.5cm,bottom=1.5cm,left=1.8cm,right=1.8cm]{geometry}
\usepackage{titlesec}
\usepackage{enumitem}
\usepackage[hidelinks,pdfusetitle]{hyperref}
\usepackage{xcolor}

% Texto selecionável/copiável e compatível com ATS
\input{glyphtounicode}
\pdfgentounicode=1

\pagestyle{empty}
\setlength{\parindent}{0pt}
\raggedbottom
\urlstyle{same}

% ---------- Seções ----------
\titleformat{\section}{\large\scshape\raggedright}{}{0em}{}[\titlerule]
\titlespacing*{\section}{0pt}{10pt}{5pt}

% ---------- Comandos ----------
\setlist[itemize]{leftmargin=1.2em,itemsep=1.5pt,topsep=2pt,parsep=0pt}

% \entry{Título}{Local}{Subtítulo}{Período}
\newcommand{\entry}[4]{%
  \vspace{3pt}
  \noindent\textbf{#1}\hfill #2\\
  \textit{#3}\hfill\textit{#4}\par
}

% \skill{Categoria}{Itens}
\newcommand{\skill}[2]{\noindent\textbf{#1:} #2\par\vspace{2pt}}


\begin{document}

% ---------- Header ----------
\begin{center}
  {\LARGE\bfseries Hellen Araujo da Silva}\\[4pt]
  São Paulo, SP \textbar{} \href{mailto:hds74062@gmail.com}{hds74062@gmail.com} \\
  \href{https://www.linkedin.com/in/hellen-araujo-da-silva-550b86322/}{linkedin.com/in/hellen-araujo-da-silva}
  \textbar{} \href{https://github.com/hellen-pyphile}{github.com/hellen-pyphile}
\end{center}

% ---------- Summary ----------
\section{Summary}
Computer Science student at PUC-SP with a technical background in Systems Development,
hands-on experience in full-stack web development and user research.
Seeking an internship in software development.

% ---------- Education ----------
\section{Education}
\entry{Pontifical Catholic University of São Paulo (PUC-SP)}{São Paulo, SP}
  {Bachelor's Degree in Computer Science}{Feb. 2026 -- Dec. 2029}
\entry{Escola Técnica Estadual Dr. Celso Giglio (ETEC)}{Osasco, SP}
  {Integrated High School with Technical Degree in Systems Development}{Feb. 2022 -- Dec. 2024}

% ---------- Projects ----------
\section{Projects}
\entry{Genpsiq -- Outreach and Academic Support in Psychiatric Genetics}{}
  {Technical Course Capstone Project (ETEC)}{2024}
\begin{itemize}
  \item \textbf{Stack:} web platform built with HTML5, CSS3 (responsive interface),
        JavaScript (client-side interactivity) and PHP (server-side business logic).
  \item \textbf{Research and design:} requirements gathered through a survey of
        52 students, of whom 76.9\% reported difficulty finding scientific
        papers; iterative prototyping in Canva to refine usability.
  \item \textbf{Features:} user authentication (sign-up and login); file management
        (PDF and DOCX uploads, ID-based indexing, filtered search and deletion);
        clinical case simulator that automatically validates anamneses and
        psychological reports.
\end{itemize}

% ---------- Experience / Activities ----------
\section{Experience and Activities}
\entry{PUCTech -- Academic League of Science and Technology}{São Paulo, SP}
  {Member (People team and project development)}{2026}

\entry{Science Fairs}{}{Presentation of the Genpsiq project}{2024}
\begin{itemize}
  \item CGTEC -- ETEC Science Fair (Osasco, SP)
  \item COTUCA -- UNICAMP Science Fair (Campinas, SP)
  \item A Journey of Knowledge and Innovation -- UNIFESP Science Fair (São Paulo, SP)
\end{itemize}

% ---------- Technical skills ----------
\section{Technical Skills}
\skill{Languages}{Python, Java, C, C++, C\#, JavaScript, SQL, HTML, CSS}
\skill{Python (libraries and frameworks)}{FastAPI, Pydantic, Requests, Kivy, IPython, json, os}
\skill{C / C++}{C Standard Library (stdio, stdlib, string); STL (vector, iostream, iomanip)}

% ---------- Languages and soft skills ----------
\section{Languages and Soft Skills}
\skill{Languages}{Portuguese (native) \textbar{} English (intermediate) \textbar{} Spanish (basic)}
\skill{Soft skills}{conflict mediation, strategic thinking, teamwork, communication and transparency}

\end{document}